\scanpage{20}

with the aid \unsure{by giving oneself} a system

\diagram{\begin{tikzcd}[column sep=large]
\text{(35)} &
\begin{gathered}
\underline{\mathcal G}\supset\underline M\\[-2pt]
\text{Group over }S_0,\\[-2pt]
\text{with }\underline M\text{ as}\\[-2pt]
\text{normal subgroup}
\end{gathered}
& \text{, and }\Phi\colon\mathcal G_{S_0}\arrow[r]
& \underline{\mathcal G}
\end{tikzcd}}

such that (34) is \unsure{reconstituted} from (35) by taking
$\underline M$, $\mathcal G_{S_0}$ in (35),
$M_{S_0}$ \deleted{\illegible{1}}$=$ the inverse image of
$\underline M\subset\underline{\mathcal G}$ under $\Phi$,
$\varphi$ induced by $\Phi$, and $\Theta$ the operation deduced from
$\Phi$. Moreover, $\Phi$ induces an \uline{isomorphism}.

\diagram{\begin{tikzcd}
\underbrace{\mathcal G_{S_0}/M_{S_0}}_{H_{S_0}}
\arrow[r,"\sim"] & \underline{\mathcal G}/\underline M
\end{tikzcd}}

-- and we have here an ``axiomatic'' description of the system (35)
in terms of (34).
The exact sequence of \deleted{\unsure{schemes}}\inserted{Groups over $\unsure{S|S_0}$}

\diagram{\begin{tikzcd}[column sep=large,row sep=small]
\text{(36)} & 1\arrow[r,dash] & \underline M\arrow[r]
& \underline{\mathcal G}\arrow[r]
& \underline{\mathcal G}/\underline M\arrow[r,dash]
\arrow[d,phantom,"\simeq"{description,rotate=90}] & 1\\
&&&& H_{S_0}.
\end{tikzcd}}

in fact, over $\forall$ affine base scheme $S$
(more generally, $S$ such that $H^1(S,M\otimes_{\mathbb Z}\underline{\mathcal O}_S)=0$)
\unsure{some} exact sequences

\margin{The hypothesis \unsure{made} about $H^1$ being \unsure{zero|zero}
is entirely useless, for the $\underline M$-torsors
\uline{over} $S_0$ defined by the \unsure{inverse} \unsure{images}
of the $h(\in H)$ are obviously trivial}

\diagram{\begin{tikzcd}[column sep=large,row sep=small]
\text{(37)} & 1\arrow[r] & \underline M(S)\arrow[r]
\arrow[d,phantom,"\simeq"{description,rotate=90}]
& \underline{\mathcal G}\unsure{(A)|(/A)}\arrow[r]
& H_{S_0}(S)\arrow[r,dash]
\arrow[d,phantom,"\simeq"{description,rotate=90}] & 1\\
&& M\otimes_{\mathbb Z}A && H\text{ if }S\text{ connected}\\[-3pt]
&&\text{if }S\text{ affine with ring }A
\end{tikzcd}}

\margin{19}

\scanpage{21}
\mspage{658}

In fact, by construction we have a commutative \unsure{diagram}
of exact sequences

\diagram{\begin{tikzcd}[column sep=large,row sep=large]
\text{(38)} & 1\arrow[r] & M_{S_0}\arrow[r]\arrow[d]
&\mathcal G_{S_0}\arrow[r]\arrow[d]
&H_{S_0}\arrow[r,dash]
\arrow[d,phantom,"\unsure{\Downarrow}"{description}] &1\\
&1\arrow[r]&\underline M\arrow[r]
&\underline{\mathcal G}\arrow[r]&H_{S_0}\arrow[r]&1
\end{tikzcd}}

which on passing to sections gives \deleted{for} $S$
\deleted{arbitrary} affine (say) connected a commutative diagram

\diagram{\begin{tikzcd}[column sep=large,row sep=large]
\text{(39)} &1\arrow[r]&M\arrow[r]\arrow[d]
&\mathcal G\arrow[r]\arrow[d]
&H\arrow[r]\arrow[d,phantom,"\unsure{\Downarrow}"{description}]&1\\
&1\arrow[r]&M\otimes_{\mathbb Z}A\arrow[r]
&\underline{\mathcal G}(A)\arrow[r,dash]&H\arrow[r,dash]&1.
\end{tikzcd}}

For $S=S_0$ we recover the announced isomorphism (28).

\margin{\uline{N.B.} As a sheaf (\unsure{without}
$S$-group structure) $\underline{\mathcal G}$ is isomorphic to the
\deleted{\unsure{scheme}}\inserted{sheaf} sum of the $H$ copies of
$\underline M$ -- it is therefore representable if $\underline M$ is.}

This construction of the $\unsure{S_0}$-group scheme
$\underline{\mathcal G}$, from an exact sequence (27), is
\unsure{manifestly} functorial \unsure{with respect|with respect} to the latter:
a \unsure{homomorphism} of exact sequences

\diagram{\begin{tikzcd}[column sep=large,row sep=large]
\text{(40)}&1\arrow[r]&M\arrow[r]\arrow[d]
&\mathcal G\arrow[r]\arrow[d]&H\arrow[r]\arrow[d]&1\\
&1\arrow[r]&M'\arrow[r]&\mathcal G'\arrow[r]&H'\arrow[r,dash]&1
\end{tikzcd}}

gives rise to a commutative diagram of exact sequences

\margin{20}

\scanpage{22}

\diagram{\begin{tikzcd}[column sep=small,row sep=large]
&1\arrow[rr,dash]&&M_{S_0}\arrow[rr]\arrow[dr]\arrow[dd]
&&\mathcal G_{S_0}\arrow[rr]\arrow[dr]\arrow[dd]
&&H_{S_0}\arrow[rr,dash]\arrow[dr]
\arrow[dd,dash,"\unsure{\Vert}"{description,pos=.25}]
&&1&\\
&&1\arrow[rr,crossing over]
&&M'_{S_0}\arrow[rr]\arrow[dd]
&&\mathcal G'_{S_0}\arrow[rr]\arrow[dd]
&&H'_{S_0}\arrow[rr]\arrow[dd]&&1\\
\text{(41)}&1\arrow[rr]&&\underline M^{\unsure{\cdot}}\arrow[rr]\arrow[dr]
&&\underline{\mathcal G}\arrow[rr]\arrow[dr]
&&H_{S_0}\arrow[rr]\arrow[dr]&&1&\deleted{\unsure{1}}\\
&&1\arrow[rr]&&\underline M'\arrow[rr]
&&\underline{\mathcal G}'\arrow[rr]
&&\underline H'\arrow[rr]&&1
\end{tikzcd}}

\margin{\uline{Discrete} \uline{schematic envelope}
\uline{of groups} \uline{that are metabelian}}

4. Let us then take up again a free $\mathbb Z$-module of finite type $M$, and for
$\forall$ sub-\deleted{\unsure{g}}\inserted{$\mathbb Z$}-module
of finite index, let us consider the exact sequence

\diagram{\begin{tikzcd}[column sep=large]
\text{(41)}&1\arrow[r,dash]&M_i\arrow[r,dash]&M\arrow[r,dash]
&M/M_i\arrow[r,dash]&1.
\end{tikzcd}}

More \unsure{generally}, if we start from an exact sequence (27), for
$\forall$ subgroup $M_i$ of finite index in $M$,
\unsure{normal} in $\mathcal G$ i.e. stable under $H$,
we have an exact sequence

\diagram{\begin{tikzcd}[column sep=large]
\text{(43)}&1\arrow[r]&M_i\arrow[r]&\mathcal G\arrow[r]
&\underbrace{\mathcal G/M_i}_{H_i}\arrow[r,dash]&1
\end{tikzcd}}

where now $H_i=\mathcal G/M_i$ \unsure{is} \unsure{the|an}
extension of

\diagram{\begin{tikzcd}[column sep=large]
\text{(44)}&1\arrow[r]&M/M_i\arrow[r]&H_i\arrow[r]&H\arrow[r]&1
\end{tikzcd}}

deduced from (27) by

\diagram{\begin{tikzcd}
M\arrow[r]&M/M_i.
\end{tikzcd}}

\margin{21}

\scanpage{23}
\mspage{659}

We note that \deleted{if $H$ is a group \unsure{having} a
finite \unsure{number} of \unsure{generators}, for $H$ finite},
\deleted{then} the $M_i\subset M$ of finite index and stable under $H$
are cofinal in the \unsure{set} of all the $M_i$ of finite index in
$M$. Let us therefore consider the associated exact sequence of group schemes

\margin{take the $nM$ !}

\diagram{\begin{tikzcd}[column sep=large]
\text{(45)}&1\arrow[r]&\underline{M_i}\arrow[r]
&\underline{\mathcal G}_i\arrow[r]&H_{i\,S_0}\arrow[r]&1
\end{tikzcd}}

For varying $i$, they form a projective system of extensions,
hence an extension of \unsure{progroups} over $S_0=\operatorname{Spec}\mathbb Z$.

\deleted{In particular}

\deleted{\uline{\unsure{Remark|The pro-group}} $\unsure{\underline{\mathcal G}}$}

We consider the pro-object $(\underline{\mathcal G}_i)$ as a pro-object
of $C$ the category of group schemes (smooth etc ...) over $S_0$.
We note that the transition morphisms

\diagram{\begin{tikzcd}
\underline{\mathcal G}_j\arrow[r]&\underline{\mathcal G}_i
\end{tikzcd}}

are \uline{affine} -- the question being local
\inserted{on $\underline{\mathcal G}_i$}, we are reduced to looking at it for
the homomorphism

\deleted{\unsure{image}}

\diagram{\begin{tikzcd}[column sep=large,row sep=small]
\unsure{E}(M_i,M_j)\arrow[r]
&E(M_i,M_i)\arrow[d,phantom,"\unsure{\Vert}"{description}]\\
M_i\supset M_j&\underline{M_i}\\[-3pt]
\text{extension of }(M_i/M_j)_{S_0}\text{ by }\underline{M_j}
\end{tikzcd}}

\margin{22}

\scanpage{24}

Now the first is affine, as a finite sum of affine schemes, the second is affine --- so the morphism is affine.

This being so, the projective limit of the $\underline{\mathcal{G}}_i$ in the category of schemes exists (more precisely, the projective limit functor of the $\underline{\mathcal{G}}_i$ is representable) --- it is perhaps not absurd to pass to the limit \unsure{outright}, and to work with this group scheme --- not of finite type.

\uline{Remark} The pro-group $(\underline{\mathcal{G}}_i)$ over $S_0=\operatorname{Spec}\unsure{\mathcal{G}|\mathbb{Z}}$ depends only on the \deleted{$\mathcal{G}$}, discrete group $\mathcal{G}$, equipped with a class of subgroups \inserted{\deleted{\unsure{invariant} abelian}} modulo commensurability --- the hypothesis to be made being that among these subgroups, there are some which are $\mathbb{Z}$-modules \deleted{free} of finite type (for taking them small enough, they are \uline{free} of finite type). If $M_0$ is itself of finite index in $\mathcal{G}$ i.e. if \deleted{the}\inserted{\unsure{one of the|$\forall$ of the}} $H_i$ are finite,

\margin{23}

\scanpage{25}
\mspage{660}

then this pro-object depends only on $\mathcal{G}$ itself
--- \unsure{subject} to the condition that there exists
a subgroup of finite index in $\mathcal{G}$ which
is a commutative group, of finite type as
a $\mathbb{Z}$-module. (We then say
that $\mathcal{G}$ is ``\uline{metabelian}''.)

\margin{\unsure{3}}
\deleted{Let us look at the typical case}

Thus, for $\forall$ discrete group
\inserted{$\mathcal{G}$} that is metabelian, we have constructed a pro-group
scheme, \inserted{projective system of \uline{affine} group schemes,}
which I will denote by
\[
(46)\qquad
\operatorname{Env}(\mathcal{G})=(\underline{\mathcal{G}}_i)
\]
(the index set being made up of the normal subgroups
of $\mathcal{G}$ which are commutative and
torsion-free). This pro-group depends \unsure{functorially}
on $\mathcal{G}$. It has properties \unsure{of exactnesses|of exactness}
that are \unsure{interesting} and $\pm$ obvious, for example
it is multiplicative
\[
(47)\qquad
\operatorname{Env}(\mathcal{G}\times\mathcal{G}')
\unsure{\xrightarrow{\sim}}
\operatorname{Env}(\mathcal{G})\mathbin{\times_{S_0}}
\operatorname{Env}(\mathcal{G}').
\]
If we have an exact sequence
\[
(48)\qquad
\diagram{\begin{tikzcd}[column sep=large]
1 \arrow[r] & \mathcal{G}'' \arrow[r,no head] &
\mathcal{G} \arrow[r,no head] & \mathcal{G}' \arrow[r,no head] & 1
\end{tikzcd}}
\]
of metabelian groups, we have an exact
sequence for the envelopes
\[
(49)\qquad
\diagram{\begin{tikzcd}[column sep=large]
1 \arrow[r] & \unsure{\operatorname{Env}}(\mathcal{G}'') \arrow[r] &
\operatorname{Env}(\mathcal{G}) \arrow[r,no head] &
\operatorname{Env}(\mathcal{G}') \arrow[r] & \unsure{1}
\end{tikzcd}}
\]

\margin{24}

\scanpage{26}

in a sense which should be made precise \unsure{on}
occasion --- for the associated group schemes
(obtained by passing to the limit) ---
\unsure{we do indeed have} a faithfully flat homomorphism whose
kernel is the right one \ldots

The typical case, to which the general case
reduces in a \unsure{fairly} immediate way, is
the one where $\mathcal{G}=M=\mathbb{Z}$. We have for
$\forall$ integer $n\geq 1$ the exact sequence
\[
(50)\qquad
\diagram{\begin{tikzcd}[column sep=large]
1 \arrow[r,no head] &
\mathbb{Z} \arrow[r,no head,"n"] &
\mathbb{Z} \arrow[r,no head] &
\mathbb{Z}/n\mathbb{Z} \arrow[r,no head] & 1
\end{tikzcd}}
\]
I am going to \deleted{\unsure{construct}} make explicit the construction
of the corresponding extension of
$\mathbb{Z}/n\mathbb{Z}$ by $\underline{\mathbb{Z}}=\mathbb{G}_a$,
i.e. $\underline{\mathcal{G}}_n$.

It is an affine scheme, whose coordinate
ring is isomorphic --- in a not
very canonical way --- to the product of the
$n$ copies of that of $\underline{\mathbb{Z}}=\mathbb{G}_a$,
namely $\mathbb{Z}[T]$ --- we find such an \unsure{isomorphism}
by \unsure{choosing} a set-theoretic section of
$\mathbb{Z}$ over $\mathbb{Z}/n\mathbb{Z}$, for example by \unsure{the set}
of integers $\mathbb{Z}\cap[0,n-1]$. The homomorphisms of

\margin{25}

\scanpage{27}
\mspage{661}

transition between the $\underline{\mathcal{G}}_n$ have nothing
flat about them of course. Passing to the projective limit,
we find an extension of
\[
(\widehat{\mathbb{Z}})_{S_0}
\overset{\mathrm{dfn}}{=}
\varprojlim_n(\mathbb{Z}/n\mathbb{Z})_{S_0}
\]
--- constant profinite group scheme over $S_0$ ---
by the projective limit of the schemes
$\mathbb{G}_a$, mapping to itself
by the homomorphisms
\unsure{seen} \deleted{\unsure{of the} \unsure{\S b}} of $2^\circ$)
--- it is the spectrum
of the inductive limit ring of the
\[
\begin{gathered}
\mathbb{Z}[T_n]\subset\mathbb{Q}[T]\\
T_n=\frac{1}{n}T
\end{gathered}
\]
which is thus, strange to say, \unsure{the} \deleted{\illegible{1}}
\inserted{\unsure{sub-}}$\mathbb{Z}$-algebra
of $\mathbb{Q}[T]$ \unsure{formed|formed} by the polynomials
for which \uline{the constant term} is an integer.

We have an analogous \unsure{description} of
the projective limit of the schematic
envelope of $\mathbb{Z}^n$ --- it is an
extension of $(\widehat{\mathbb{Z}}^{\,n})_{S_0}$
(beware of the abuse
of notation) by an identity component,
which is the spectrum of the subring
of $\mathbb{Q}[T_1,\ldots,T_n]$ formed by the polynomials whose
constant term is an integer \ldots

\margin{26}